
\section{Experiment}

The data presented here were taken using the \minerva detector and the wideband antineutrino beam produced by the NuMI beamline
in the low-energy mode~\cite{Anderson:1998zza} with a mean energy of $3.6$~GeV.
The antineutrino flux is estimated from a simulation of the neutrino beamline based on Geant4~\cite{Agostinelli2003250,1610988}, with
hadron production in the simulation constrained by proton-carbon external data~\cite{Alt:2006fr,Barton:1982dg,Lebedev:2007zz}.
The \minerva detector consists of a fully active region of scintillator strips surrounded on the sides 
and downstream end by electromagnetic and hadronic calorimeters, and is described in detail in Ref.~\cite{minerva_nim}.
There are three orientations (views) for the strips (X,U,V), offset by $60^\circ$ from each other, which enable 
three-dimensional reconstruction of particle trajectories.
The X view is sampled twice as often as the other views.
The downstream edge of the \minerva detector is located \unit[2]{m} upstream of the MINOS Near Detector, 
a magnetized iron spectrometer~\cite{Michael:2008bc} used in this analysis 
to reconstruct the momentum and charge of muons.
The transport of particles from neutrino interactions in the detector is simulated by a Geant4-based program.
The readout simulation is tuned so that both the photostatistics and the reconstructed energy deposited
by momentum-analyzed through-going muons agree between data and the simulation. The detector 
simulation of single particle responses is validated using testbeam data taken with a scaled-down 
version of the MINERvA detector~\cite{Aliaga:2015aqe}.

Neutrino interactions are simulated using the GENIE 2.6.2 neutrino event generator. 
Details concerning GENIE and its associated parameters are described in Ref.~\cite{Andreopoulos201087}.
For baryon resonance production, the formalism of Rein-Sehgal~\cite{Rein:1980wg} is used with modern resonance properties~\cite{pdg}.
Non-resonant pion production is simulated using the Bodek-Yang model~\cite {Bodek:2004pc} and is constrained below $W=1.7$~GeV 
by neutrino-deuterium bubble chamber data~\cite{Radecky:1981fn,Kitagaki:1990vs}.
Pion and nucleon FSI are modeled in GENIE using a parameterized intranuclear cascade model, 
with the full cascade being represented by a single interaction.  For all models fitting data for light nuclei such as carbon, 
the pion most often has one interaction as it propagates through the nucleus.  
For each interaction, choice of the channel (e.g. charge exchange) is based on
total cross section data and calculations~\cite{Lee:2002eq,Ashery:1981tq} and the kinematics and multiplicity of the final state are taken from fits to more detailed data.
At the energies important for this measurement, the hadron-nucleus cross sections have large uncertainties.
For example, the cross section of $\pi^-\rightarrow \pi^0$ CEX on carbon has an uncertainty of about 50\%.
Reaction cross sections for $\pi^0$ are estimated from measured cross sections for $\pi^\pm$ scattering using isospin symmetry or theoretical models.
The most significant advantages of the single interaction used in GENIE's FSI model
are the ability to exactly reweight and to characterize each event with a single final state channel.
The uncertainties in the FSI model are evaluated by varying its strength within previously 
measured hadron-nucleus cross-section uncertainties.

This analysis uses data taken between October 2009 and February 2012 with $2.01\times 10^{20}$ protons on target (POT) in the $\bar{\nu}_\mu$ mode. 
About half of the exposure ($0.945\times 10^{20}$ POT) was taken during construction with the downstream half of the detector.
In this period the ArgoNeuT detector~\cite{Argoneut} was situated between the \minerva and MINOS Near detectors. 
Because the two sub-samples of the data have different efficiencies, they were analyzed separately and their results combined.

